% Keep styling.sty beside this file. Compile with:
% latexmk -xelatex -synctex=1 template.tex (or, 
% for my purposes, compile to preview with VimtexCompile in neovim)
% if youre using VSC with some built in plugin you'll probably need to add 
% {
%   "latex-workshop.latex.recipe.default": "latexmk (xelatex)",
%   "latex-workshop.latex.autoBuild.run": "onSave",
%   "latex-workshop.view.pdf.viewer": "tab",
%   "latex-workshop.view.pdf.invertMode.enabled": "never"
% }
% also it is necessary to download the right fonts. Just look 
% at error messages to find which ones or something

\documentclass[11pt]{article}
\usepackage{styling} % Dark by default, with the clickable THEME switch.
% \usepackage{graphicx} % Uncomment if adding external images.

\usepackage{algpseudocode}

\title{Algorithms assignment 2}
\author{Thobias K. H\o ivik}
\date{For November 13, 2026} % Use \date{} to hide the date.

\begin{document}
\maketitle


\begin{prob*}[2 Problem 5]
    Given $7$ jobs with corresponding 
    processing times, 
    \begin{center}
        \begin{tabular}{ |c|c|c|c|c|c|c|c| } 
            \hline
            Job $j$ & 1 & 2 & 3 & 4 & 5 & 6 & 7 \\ 
            Processing Time $p_j$ & 4 & 7 & 2 & 9 & 12 & 5 & 3 \\ 
            \hline
        \end{tabular}
    \end{center}

    The jobs will be processed on $3$ identical machines 
    without interruption. All jobs are avliable at 
    time $0$. (Model on slide 13, Chapter 2)
    
    \begin{enumerate}[label=\alph*)]
        \item Use the local search algorithm where we have 
            all jobs on machine $1$ in the beginning (in the 
            same order as they are given). When you move 
            a job, move it to the machine it can start 
            earliest. I there is a tie, use the machine with the lowest number. What is the result, both
            time when the last job finish and the actual schedule?
        \item Use the Greedy algorithm on the same data. When have all jobs finished?
        \item Use the Longest processing time rule. When have all jobs finished?
        \item Give a lower bound on the best solution and give a short explanation?
        \item Is it possible to achieve the answer in d) (Explain)?
    \end{enumerate}
\end{prob*}

This is the classical identical parallel-machine makespan 
minimization problem, denoted 
\[
    P_3 \| \ \| C_{\max}. 
\]
We want to find a schedule minimizing 
\[
    C_{\max} = \max_{i\in[3]} L_i,  
\]
where 
\[
    L_i := \sum_{j \ \text{assigned to machine} \ i} p_j
\]
is the total processing time on machine $i$.


\begin{proof}[Solution to (a)]
    The initial schedule gives all seven jobs to 
    machine $1$ in the described order by 
    \[
        M_1 = \underbrace{1}_{4}
        \qquad 
        \underbrace{1}_{4}
        \qquad 
        \underbrace{2}_{7}
        \qquad 
        \underbrace{3}_{2}
        \qquad 
        \underbrace{4}_{9}
        \qquad 
        \underbrace{5}_{12}
        \qquad 
        \underbrace{6}_{5}
        \qquad 
        \underbrace{7}_{3}.
    \]
    Note: This accidentally looks a lot like fractions, but 
    I am just illustrating the processing times $p_j$. 

    We get initial loads 
    \[
        (L_1,L_2,L_3) = (42,0,0)
    \]
    with $C_{\max} = 42$. 

    If $M_j$ is the unique machine for which 
    $L_j = C_{\max}$ we find that it is advantageous 
    to swap if 
    \[
        L_j + c < L_i, 
    \]
    or equivalently 
    \[
        c < L_i - L_j.
    \]

    We perform the following procedure. 

    \begin{algorithmic}[1]
        \Require Jobs $J$ with processing times $p_j$,
                 machines $M_1,\ldots,M_m$
        \Ensure A locally optimal schedule

        \State Assign all jobs to $M_1$ in their given order
        \State $L_i \gets \sum_{j \in M_i} p_j$
               for each $i \in \{1,\ldots,m\}$

        \While{true}
            \State $i \gets \arg\max_{k} L_k$
            \State $\ell \gets$ last job on $M_i$
            \State $j \gets \arg\min_{k \neq i} L_k$

            \If{$L_j + p_\ell \geq L_i$}
                \State \Return $(M_1,\ldots,M_m)$
            \EndIf

            \State Remove $\ell$ from $M_i$
            \State Append $\ell$ to $M_j$
            \State $L_i \gets L_i - p_\ell$
            \State $L_j \gets L_j + p_\ell$
        \EndWhile
    \end{algorithmic}

    So here we begin with 
    \[
        (42,0,0). 
    \]
    The job which finishes last is $j = 7$ so we get 
    \[
        (42,0,0) \mapsto (39,3,0). 
    \]
    Now $M_1$ still is the machine with maximum load 
    so we take $j = 6$, and move 
    \[
        (39,3,0) \mapsto (34, 3, 5). 
    \]
    Repeat to get 
    \[
        (34,3,5) \mapsto (22,15,5) 
        \mapsto (15,15,14). 
    \]
    Now we look at $M_1$ again and see that 
    for $j = 3$ we have $p_j = 3 < L_3 = 14$ so 
    the algorithm halts with 
    \[
        C_{\max} = 14 \leq 2C^*_{\max} = 
        14 \cdot 2 = 28
    \]
    as expected. 

    Running time is $O(nm)$.

\end{proof}

\begin{proof}[Solution to (b)]
    We perform 
    \begin{algorithmic}[1]
        \Require Jobs with processing times
                 $p_1,\ldots,p_n$ and $m$ machines
        \State $L_i \gets 0$ for all $i=1,\ldots,m$
        \For{$j=1,\ldots,n$}
            \State $i^* \gets \arg\min_i L_i$
            \State Assign job $j$ to machine $M_{i^*}$
            \State $L_{i^*} \gets L_{i^*}+p_j$
        \EndFor
        \State \Return $\max_i L_i$
    \end{algorithmic}

    We begin at $j = 1$ and since all $M_i$ have 
    $L_i = 0$ we get $L_i$'s 
    \[
        (4,0,0). 
    \]
    Next 
    \[
        (4,7,0), 
    \]
    then 
    \[
        (4,7,2).
    \]
    Once again $L_3$ is the minimum so next pass gives 
    \[
        (4,7,11). 
    \]
    Now $L_1$ is smallest so we get 
    \[
        (16,7,11). 
    \]
    Now $L_2$ is smallest so we get 
    \[
        (16,12,11). 
    \]
    Now $L_3$ is smallest so we get 
    \[
        (16,12,14). 
    \]
    We obtain $C_{\max} = 16$ with 
    the algorithm running in time $O(nm)$. 
\end{proof}

\begin{proof}[Solution to (c)]
    We perform 
    \begin{algorithmic}[1]
        \Require Jobs $J$ with processing times $p_j$,
                 $m$ identical machines
        \State Sort jobs by nonincreasing processing time
        \State $L_i \gets 0$ for all $i=1,\ldots,m$
        \For{each job $j$ in sorted order}
            \State $i^* \gets \arg\min_i L_i$
            \State Assign job $j$ to machine $M_{i^*}$
            \State $L_{i^*} \gets L_{i^*}+p_j$
        \EndFor
        \State \Return $\max_i L_i$
    \end{algorithmic}

    Performiing the algorithm we get 
    \begin{align*}
        (12,0,0) &\mapsto (12,9,0) \\ 
        &\mapsto (12,9,7) \\ 
        &\mapsto (12,9,12) \\ 
        &\mapsto (12,13,12) \\ 
        &\mapsto (15,13,12) \\ 
        &\mapsto (15,13,12)
    \end{align*}
    attaining a maximum of $15$. 

    Running time is 
    $O(n\log n + mn)$ since we sort first. 
\end{proof}

\begin{proof}[Solutioin and proof sketch for (d)]
    Let $C_{\max}^*$ denote the optimal makespan, and 
    let $L_i$ be the total processing time assigned 
    to machine $M_i$. 

    Since every job must be processed, we have 
    \[
        \sum_{i=1}^m L_i = \sum_{j=1}^n p_j. 
    \]
    Moreover, since each machine processes jobs 
    sequentially, 
    \[
        C_{\max}^* \geq \max_i L_i \geq \frac 1 m 
        \sum_{i=1}^m L_i = \frac 1 m \sum_{j=1}^n p_j. 
    \]
    We also necessarily have 
    \[
        C_{\max}^* \geq \max_j p_j, 
    \]
    since every job must finish before the 
    schedule is complete. 

    Combining we get 
    \[
        C_{\max}^* \geq 
        \max \left\{
            \frac 1 m \sum_{j=1}^m p_j, \ \max_j p_j
        \right\}.
    \]
\end{proof}

\begin{proof}[Solution to (e)]
    Consider the following schedule, 
    \begin{align*}
        M_1 : \{5,3\},  & L_1 = 12 + 2 = 14 \\ 
        M_2 : \{4,6\},  & L_2 = 9+5 = 14, \\ 
        M_3 : \{2,1,7\} & L_3 = 7+4+3 = 14.
    \end{align*}

    This is a feasible schedule with 
    \[
        C_{\max} = \max\{14,14,14\} = 14 = C_{\max}^*. 
    \]
\end{proof}

\begin{prob*}[2, Problem 6]
    The adjacency matrix of a graph on $6$ vertices 
    is given below. The triangle inequality holds for the 
    whole graph. 
    \[
        \begin{bmatrix}
            0 & 4 & 5 &10&11& 7 \\ 
            4 & 0 & 9 & 6& 7& 9 \\ 
            5 & 9 & 0 & 9& 6& 2 \\ 
           10 & 6 & 9 & 0& 3& 9 \\ 
           11 & 7 & 6 & 3& 0& 6 \\ 
            7 & 9 & 2 & 9& 6& 0 
        \end{bmatrix}
    \]
    \begin{enumerate}[label=\alph*)]
        \item Verify that the triangle inquality 
            holds when going from vertex $1$ to $5$. 
        \item Run the nearest addition algorithm on 
            the graph. If there are more than 
            one candidate to add in a step, use the vertex 
            with the smallest number. It is sufficient to 
            have fours nodes in the tour (skipping the 
            last two iterations through the loop).
        \item Run the Double tree algorithm on the graph. 
            If there are more than one minimum spanning 
            tree, you should run it on every tree. 
        \item Run Christofides' algorithm on the graph. 
            If there are more than one minimum spanning 
            tree, you should run it on every tree. 
    \end{enumerate}
\end{prob*}

Let $G = ([6], E) \cong K_6$ be the complete 
graph (undirected) on $6$ vertices with edge weights given 
in the adjacency matrix in the problem description.

Denote the weights $c_{ij}$. I.e. 
$c_{12} = 4, \ c_{15} = 11$ and so forth. 

\begin{proof}[Solution for (a)]
    We have $c_{15} = 11$. 
    The possibilities are 
    \[
        \begin{aligned}
        c_{12}+c_{25}&=4+7=11,\\
        c_{13}+c_{35}&=5+6=11,\\
        c_{14}+c_{45}&=10+3=13,\\
        c_{16}+c_{65}&=7+6=13.
        \end{aligned}
    \]
    Thus 
    \[
        c_{15} = 11 \leq \min\{11,11,13,13\} = 11 
    \]
    so were good.
\end{proof}

\begin{proof}[Solution to (b)]
    We wish to find a permutation $\pi$ of the 
    vertices minimizing 
    \[
        c_{\pi(1)\pi(2)} + 
        c_{\pi(2)\pi(3)} + \cdots + 
        c_{\pi(6)\pi(1)}. 
    \]

    I.e. we want to find 
    a minimum-weight Hamiltonian cycle. 
    The additional assumption is that the weights 
    satisfy the triangle inequality 
    \[
        c_{ij} \leq c_{ik} + c_{ki}, \ \forall i,j,k \in V = [6].
    \]

    We begin with the two closest vertices 
    and repeatedly add the vertex outside 
    the tour that is closest to any vertex 
    already in the tour. 

    When inserting a new vertex $j$ between 
    two consecutive vertices $u,v$, the increase in 
    tour cose is 
    \[
        \Delta = c_{uj} + c_{jv} - c_{uv}.
    \]
\end{proof}

\end{document}
